\chapter{重整化、自然性与层级问题}

\lectureoverview{本课程的核心思想源于标准模型中关于重整化的种种疑难。因此，我们首先讨论重整化的物理含义：消去对于当前问题而言过小、过快而无关的自由度，然后将其效应编码进一个有效描述之中。一个分子提供了原型。随后，量纲分析将同样的推理带入标量量子场论，在那里它揭示出标量质量的二次敏感性、希格斯场的层级问题、费米子与规范玻色子所享有的保护，以及更为严重的真空能问题。}

\section{问题尺度上的物理}

本季度将要展开的大部分思想，都起源于标准模型中关于重整化的问题---有时是疑难，有时是有用的观察---。重整化结合了两种思维习惯。第一种是从描述中去除那些距离短到与所问问题无关的结构。第二种是量纲分析，它常常在计算任何困难积分之前就已确定短距离下的答案。

这种对细节的消去贯穿整个物理学。如果我们想理解一个原子核，可以从夸克和量子色动力学出发。在计算出质子和中子的质量与自旋，以及它们的有效核力之后，对于许多目的而言，我们可以忘掉它们的夸克子结构。核运动足够缓慢，以至于质子和中子的非相对论量子力学往往是更有用的描述。

在下一个层次上，原子物理学通常不需要原子核的内部构成。核物理学告诉我们存在哪些原子核和同位素，并提供诸如质量和电荷之类的性质。于是我们以一组新的有效基本对象重新开始：原子核和电子。夸克、质子和中子并没有不复存在；它们的效应已被压缩进原子核所携带的参数之中。

保留多少取决于下一个问题。原子计算可能需要核质量，因为它进入约化质量；它当然需要电荷，因为电荷决定库仑场。它通常不需要知道哪个夸克携带核动量的哪一部分。元素电荷大约有九十种，若计入同位素则核种类更多。这份目录可以作为输入。另一方面，电子还不能被消去：电子的运动恰恰是原子结构的成因。

在解决了原子问题之后，我们同样可以把原子视为分子的构建单元。Susskind 那异想天开的清单——曲奇原子、咖啡原子、巧克力原子和番茄酱原子——说明了一个严肃的观点：每一个层次都开启一种新的描述。一旦分子性质和分子力已知，分子就可以被拼接成更大的、类似 Tinkertoy 的组装体，而不必在每一步重新打开电子问题。

每一次过渡都是一次粗粒化。新的描述在绝对意义上可能不那么精确，但在其预期的尺度上却有用得多。量子场论使这个问题无法回避，因为场在每一种波长上都有模式，包括任意短的波长。因此，我们把短到无法分辨的波长的效应求和，并用新的有效参数取而代之。这一操作，连同量纲分析，就是重整化的物理内容。

\section{作为原型的分子}

考虑两个原子，并寻求对它们相对运动的简单描述。每个原子都包含一个沉重的原子核和一团电子云。即便质子也大约比电子重两千倍，而更重的原子核会进一步拉大尺度之间的间隔。一个有用的图像是：一对保龄球被快速运动的苍蝇所环绕——原子核响应缓慢，而电子态则迅速重新排布。

设原子核的位置和动量分别为 \(R_a\) 和 \(P_a\)，共同质量为 \(M\)；设电子的坐标和动量为 \(r_i\) 和 \(q_i\)，质量为 \(m_e\)。略去诸如 \(4\pi\epsilon_0\) 之类的常规因子后，哈密顿量具有如下示意形式
\begin{align}
H={}&
\frac{P_1^2}{2M}+\frac{P_2^2}{2M}
+\frac{e^2}{R_{12}}
+\sum_i\frac{q_i^2}{2m_e}
+\frac12\sum_{i\ne j}\frac{e^2}{r_{ij}}
\nonumber\\
&-\sum_i\left(
\frac{e^2}{|R_1-r_i|}
+\frac{e^2}{|R_2-r_i|}
\right).
\label{eq:molecular-hamiltonian}
\end{align}
第一行包含原子核的动能、它们之间的库仑排斥以及电子的动能。第二行包括电子-电子排斥和电子-原子核吸引。对于具有一般电荷的原子核，相应的 \(Z_a\) 因子应予恢复；本讲之所以将电荷保持为示意形式，是因为这里的关键问题是时间尺度的分离。

有两种特征运动。原子核重而慢；电子轻而快，快到看起来像一团模糊。在第一近似中，固定核的位置，并把所有依赖于电子的项归入一个电子哈密顿量，
\begin{equation}
H_{\mathrm{el}}(R_1,R_2).
\end{equation}
对每一个固定的核构型，求解
\begin{equation}
H_{\mathrm{el}}(R_1,R_2)
\psi_0(r_1,\ldots,r_N;R_1,R_2)
=
E_{\mathrm{el}}(R_1,R_2)\psi_0.
\label{eq:electronic-ground-state}
\end{equation}
我们保留最低的电子本征值。电子并没有被忽略；它们所有的快速运动都已由函数
\(E_{\mathrm{el}}(R_1,R_2)\) 概括。

其逻辑是绝热的。在一次快速的电子调整期间，原子核移动得如此之少，以至于它们的坐标可以被当作参数。在电子本征问题求解之后，又允许这些相同的坐标再次移动。在每一个缓慢变化的构型下，电子云都跟随其瞬时基态。更精确的处理可以保留激发态电子表面以及它们之间的跃迁，但领头表面已经体现了重整化思想。

\begin{figure}[htbp]
\centering
\includegraphics[width=0.72\textwidth]{lecture_01_figure_02.png}
\caption{在引入电子基态问题时，分子哈密顿量仍留在黑板上。式~\eqref{eq:molecular-hamiltonian}中的排版形式记录了此画面中被遮挡的项。}
\lectureframe{00:16:19}
\end{figure}

\begin{classroomqa}
\audiencequestion{核动量只是被简单地设为零吗？\lecturetimestamp{00:15:39}}
\lecturerresponse{只是暂时如此。在领头近似中，原子核被当作钉住不动，因为它们比电子慢得多。当我们求解有效核问题时，它们的动量将被恢复，因此这一近似并不违反不确定性原理。}
\end{classroomqa}

电子基态能量现在是原子核所感受到的势的一个贡献：
\begin{equation}
V_{\mathrm{eff}}(R_1,R_2)
=\frac{e^2}{R_{12}}+E_{\mathrm{el}}(R_1,R_2).
\label{eq:molecular-effective-potential}
\end{equation}
转动和平移不变性将它对于两个孤立原子约化为其间距 \(R_{12}\) 的函数。在非常短的距离上，核库仑排斥占主导。在非常大的距离上，相互作用很小。在这些极限之间，电子云可以产生一个吸引区域和一个极小值。于是核波函数就位于该极小值附近，并描述一个束缚分子。

注意运算的顺序。我们并没有猜测一个经验吸引力然后丢弃电子。我们首先在固定的核背景下求解了它们的量子问题。只有在那之后，我们才抹去它们的显式坐标并保留所得到的本征值。短距离物理存留在 \(V_{\mathrm{eff}}\) 的形状、深度和曲率之中。其极小值的位置决定了平衡间距，其曲率控制小振动，而其渐近形式控制束缚态是否存在。一个有效函数把所有这些电子信息带入较慢的问题中。

因此，一个复杂的多电子哈密顿量已经变成了一个带有新势的两体核哈密顿量。我们消去了快自由度，但保留了它们完整的低能效应。该势已被重整化。这一熟悉的 Born--Oppenheimer 推理通常并不以这个名称被引入，然而它恰恰包含了我们需要的思想：
\begin{equation}
\begin{gathered}
\text{fast microscopic dynamics}\\
\Downarrow\\
\text{effective parameters for slow dynamics}
\end{gathered}
\end{equation}

\section{自然单位与场的量纲}

在量子场论中，要被消去的模式可以等价地描述为短距离、短波长、高频或高能模式。在画图之前，让我们先确定量纲记账。

普通力学使用长度、时间和质量的独立单位。在粒子物理学中，我们设定
\begin{equation}
\hbar=c=1.
\end{equation}
于是能量和质量通过 \(E=mc^2\) 等价，长度和时间等价是因为 \(c=1\)，而 \(\hbar=1\) 使动量成为长度的倒数。只剩下一个量纲：
\begin{equation}
[M]=[E]=[p]=L^{-1},
\qquad
[x]=[t]=L.
\label{eq:natural-units}
\end{equation}

这并没有抹去量纲。它选择换算因子，使得时间可以被引用为长度，质量可以被引用为长度的倒数。一旦指定了一个单位，其他每一个单位都随之确定。这就是为什么截断既可以描述为短距离 \(\delta\)，也可以描述为高能量 \(\Lambda=1/\delta\)：它们是同一种尺度的两种语言。

取一个实标量场 \(\phi\)，其示意性归一化的拉格朗日密度为
\begin{equation}
\mathcal L=(\partial_\mu\phi)(\partial^\mu\phi)-V(\phi),
\qquad
S=\int d^4x\,\mathcal L.
\label{eq:scalar-action}
\end{equation}

\begin{figure}[htbp]
\centering
\includegraphics[width=0.78\textwidth]{lecture_01_figure_03.png}
\caption{局域标量拉格朗日密度及其在时空上的积分，即作用量。}
\lectureframe{00:26:08}
\end{figure}

拉格朗日密度是局域的：它描述一个时空点。经典场方程通过使作用量取驻值而得到。在量子理论中，同一个作用量出现在路径积分相位 \(\exp(iS)\) 中。无论从哪种观点看，\(S\) 的量纲都是决定性的。恢复 \(\hbar\) 后，相位为 \(\exp(iS/\hbar)\)；在设定 \(\hbar=1\) 之后，\(S\) 本身必须是无量纲的。

\begin{classroomqa}
\audiencequestion{作用量的单位是\(\hbar\)吗？\lecturetimestamp{00:26:13}}
\lecturerresponse{是的。既然\(\hbar\)已被取为一，作用量在这些单位下就是无量纲的。四维测度的量纲为\(L^4\)，因此拉格朗日密度必须具有量纲\(L^{-4}\)。}
\end{classroomqa}

动能项现在确定了场的量纲：
\begin{align}
[S]&=1, & [\mathcal L]&=L^{-4},\\
[\partial_\mu]&=L^{-1}, &
[(\partial\phi)^2]&=L^{-2}[\phi]^2=L^{-4},\nonumber\\
&& [\phi]&=L^{-1}.
\label{eq:scalar-dimension}
\end{align}
这个小小的计算将决定随后每一个图在紫外区的量纲。

将势展开为
\begin{equation}
V(\phi)=\frac{m^2}{2}\phi^2+g\phi^3+\lambda\phi^4+\cdots.
\label{eq:scalar-potential}
\end{equation}
每一项都必须具有量纲\(L^{-4}\)。由于\([\phi^n]=L^{-n}\)，
\begin{equation}
[m^2]=L^{-2},
\qquad
[g]=L^{-1},
\qquad
[\lambda]=1.
\label{eq:coupling-dimensions}
\end{equation}
四次耦合是无量纲的。这类耦合在重整化中尤为重要，因为它们的标度依赖性无法从长度的简单幂次中读出。

\begin{classroomqa}
\audiencequestion{\(g\)具有质量的单位吗？\lecturetimestamp{00:29:52}}
\lecturerresponse{是的。因子\(\phi^3\)的量纲为\(L^{-3}\)，而势密度必须具有量纲\(L^{-4}\)。因此\([g]=L^{-1}\)，在自然单位制中这是一个质量。}
\end{classroomqa}

\section{顶点与短距离传播子}

费曼规则由顶点和传播子构成。\(g\phi^3\)项给出一个三线顶点，携带一个正比于\(g\)的因子；\(\lambda\phi^4\)项给出一个四线顶点，携带一个正比于\(\lambda\)的因子。二次项给出一个局域的两点插入。它可以被想象为一个粒子进入、被吸收、然后被重新发射，其系数正比于\(m^2/2\)。数值归一化、符号、\(i\)的因子以及对称因子都被有意省略了，因为本讲要提取的是标度依赖性，而不是计算截面。

这些图形有一种操作性的解读。一个三点顶点可以吸收两个量子并发射一个，或者代表其三条线的任何与交叉相关的分配。一个四点顶点可以描述两个粒子进入、两个粒子离开。二次交叉仍然是一个合法的两线顶点，尽管它通常被并入自由传播子中。我们展示它，是因为每一个未解析的两点过程很快都将通过与它比较而被辨认出来。

传播子连接两个时空点。它的线并不表示一个量子粒子沿着一条经典的直线路径运动。它代表在\(x\)处产生一个场激发、并在\(y\)处将其移除或探测到的振幅：
\begin{equation}
\Delta(x-y)=\langle0|\phi(y)\phi(x)|0\rangle.
\label{eq:position-propagator}
\end{equation}
从右向左读这个表达式。用\(\phi(x)\)作用于真空，在\(x\)处产生一个单粒子分量。用\(\phi(y)\)作用于左矢，移除\(y\)处的分量。它们的真空矩阵元就是第一点处产生与第二点处探测之间的重叠振幅。符号\(x\)和\(y\)表示整个四维矢量，而不仅仅是两个笛卡尔坐标。

态不携带任何物理量纲，而每个场贡献\(L^{-1}\)。因此
\begin{equation}
[\Delta]=L^{-2}.
\end{equation}
对于无质量场，不存在内禀长度标度，因此量纲分析迫使短距离行为为
\begin{equation}
\Delta(x-y)\sim\frac{1}{(x-y)^2}.
\label{eq:massless-propagator}
\end{equation}
分母是洛伦兹不变的间隔平方，在其奇点处需要通常的规定。\editorialnote{对于幂次计数，方便的做法是威克转动到欧几里得间隔，在那里其大小明确地正比于\(1/|x-y|^2\)。这不会改变下文所用的紫外幂次。}

当\(y\)趋近于\(x\)时，传播子发散。这种在间隔趋于零时传播的发散，正是更复杂图中紫外发散得以建立的源头。

\section{标量质量如何被重整化}

对于标量场，动力学参数是\(m^2\)，正如相对论色散关系已经暗示的那样：
\begin{equation}
E=\sqrt{\mathbf p^2+m^2}.
\end{equation}
设\(\delta\)为我们的有效描述所能分辨的最短距离。我们不会区分顶点间距小于\(\delta\)的过程。任何这样的未解析过程，只要有一个入射和一个出射标量，就具有与质量插入相同的可观测两点形状，因此必须对有效质量有贡献。

把加速器想象成一台显微镜。如果它最精细的分辨率是\(\delta\)，那么在一个点处吸收、随后在距离远小于\(\delta\)处发射，就无法与一个真正局域的事件区分开来。我们有意在那个标度上模糊顶点。重整化就是这样一条规则：把所有这类低于分辨率的替代方案都转移到相应局域算符的系数中去。

\subsection{第一个四次修正}

从一个单独的\(\phi^4\)顶点开始。两条腿是外线粒子；另外两条腿结合成一个短圈，在几乎同一点被发射并重新吸收。将重合的传播子在一个截断距离上抹平，得到
\begin{equation}
\Delta(0)\ \longrightarrow\ \Delta(\delta)\sim\frac{1}{\delta^2},
\end{equation}
因此
\begin{equation}
\delta m^2_{(1)}\sim\frac{\lambda}{\delta^2}.
\label{eq:one-loop-scalar-mass}
\end{equation}
如果\(m_0^2\)表示输入参数或裸参数，那么低分辨率两点系数的示意形式为
\begin{equation}
\frac{m_{\mathrm{eff}}^2}{2}
\sim
\frac{m_0^2}{2}
+c_1\frac{\lambda}{\delta^2},
\label{eq:first-renormalized-mass}
\end{equation}
其中\(c_1\)包含被忽略的数值因子和约定依赖因子。在无法分辨小于\(\delta\)的距离的实验中测得的实验室质量是这个重整化后的量，而不是单独的\(m_0\)。

因此，所观测到的两点事件有两种不可区分的方式发生。原来的二次顶点可以吸收并重新发射粒子，或者四次相互作用可以做到这一点，同时短暂地发射并重新吸收一个虚标量。量子力学将振幅相加。一旦圈位于分辨率尺度以下，它的独立历史就从有效描述中消失了，但它对被测系数的贡献并没有消失。

\subsection{下一个四次修正}

现在连接两个四次顶点，使得一个外线粒子从\(x\)进入，一个从\(x+\Delta\)离开，并且三条内部传播子连接这两个顶点。有\(\lambda\)的两个幂次，而传播子给出\(\Delta^{-6}\)。量子振幅必须对出射粒子可能出现的每一个未分辨位置求和，因此对相对四维矢量进行积分：
\begin{equation}
\delta m^2_{(2)}
\sim
\lambda^2
\int_{|\Delta|\gtrsim\delta}
d^4\Delta\,\frac{1}{|\Delta|^6}.
\label{eq:two-loop-power-count}
\end{equation}

\begin{classroomqa}
\audiencequestion{顶点处的小标记是否意味着那里只有三个粒子？\lecturetimestamp{00:46:35}}
\lecturerresponse{不。我们仍然在使用四次相互作用。每个顶点有四条线相交；在这个两点图中，一条线是外线，三条线在两个顶点之间延伸。第二个顶点同样有四条入射线。}
\end{classroomqa}

测度提供长度的四个幂次，而三个传播子提供负六。由于紫外截断是估计中唯一的尺度，
\begin{equation}
\int_{|\Delta|\gtrsim\delta}
\frac{d^4\Delta}{|\Delta|^6}
\sim\frac{1}{\delta^2},
\qquad
\delta m^2_{(2)}
\sim\frac{\lambda^2}{\delta^2}.
\label{eq:two-loop-scalar-mass}
\end{equation}
这里\(\Delta^\mu\)是一个四分量积分变量；对其时间分量和三个空间分量进行积分。小写\(\delta\)是固定的截断。讲座最初对两者使用相同的口头名称，后来更正了记号。

通过将四维测度写成径向形式，可以看到同样的结果。忽略一个角度常数，
\begin{equation}
\int_{\delta}\frac{d^4\Delta}{|\Delta|^6}
\sim
\int_{\delta}\frac{R^3\,dR}{R^6}
=\int_{\delta}\frac{dR}{R^3}
\sim\frac{1}{\delta^2}.
\end{equation}
这就是为什么大距离上限对紫外估计不重要：积分由其下限端点主导。

更高阶的图重复这一模式。具有更多四次顶点的图会增加\(\lambda\)的更高幂次，但量纲分析仍然允许二次发散的質量贡献：
\begin{equation}
\begin{aligned}
m_{\mathrm{eff}}^2
={}&m_0^2
+\frac{1}{\delta^2}
\left(c_1\lambda+c_2\lambda^2+c_3\lambda^3+\cdots\right)\\
&+\text{weaker terms}.
\end{aligned}
\label{eq:scalar-mass-series}
\end{equation}

\subsection{三次相互作用与一个对数}

两个三次顶点在邻近点之间包含两个传播子。它们的贡献按如下方式标度：
\begin{equation}
\delta m_g^2
\sim
g^2\int_{|\Delta|\gtrsim\delta}
\frac{d^4\Delta}{|\Delta|^4}.
\label{eq:cubic-loop}
\end{equation}
现在测度和分母带有相同的幂次。剩余的依赖是对数性的，而不是幂次性的：
\begin{equation}
\delta m_g^2
\sim
c_g g^2\log\!\left(\frac{1}{\mu\delta}\right).
\label{eq:cubic-log}
\end{equation}
参考尺度\(\mu\)使对数无量纲。这是讲座中示意性的\(g^2\log\delta\)的完全量纲版本。对数比任何逆幂变化得慢得多，因此当截断变短时，\(1/\delta^2\)占主导。

径向形式再次使这一区别变得透明：
\begin{equation}
\int_{\delta}^{R_{\mathrm{IR}}}
\frac{d^4\Delta}{|\Delta|^4}
\sim
\int_{\delta}^{R_{\mathrm{IR}}}\frac{R^3\,dR}{R^4}
=\log\!\left(\frac{R_{\mathrm{IR}}}{\delta}\right).
\end{equation}
测度和分母中相同的幂次产生对数。分母中更多的幂次产生截断的逆幂。讲座并不是要求我们计算角度因子；它是在教我们一眼认出紫外答案。

\begin{classroomqa}
\audiencequestion{对数取什么底？\lecturetimestamp{00:57:42}}
\lecturerresponse{使用自然对数。改变底只会将它乘以一个数值常数，而该常数被吸收到系数中，其\(2\)、\(\pi\)和对称性数的因子在此尚未计算。}
\end{classroomqa}

\begin{classroomqa}
\audiencequestion{如何对一个长度取对数，质量修正的单位又来自哪里？\lecturetimestamp{00:58:32}}
\lecturerresponse{对数总是比较两个尺度；严格地说，它是像\(1/(\mu\delta)\)这样的无量纲比值的对数。耦合\(g\)具有一个质量幂次，因此\(g^2\)已经恰好具有修正所需的质量平方量纲。}
\end{classroomqa}

\begin{classroomqa}
\audiencequestion{费曼图给出一个振幅。它的模平方是概率吗？\lecturetimestamp{00:59:41}}
\lecturerresponse{是的，当该振幅被置于一个物理跃迁计算中，并带有必要的态因子和相空间因子时，它就是概率。而在这里，它更直接的作用有所不同：一个未解析的两点振幅与质量项具有相同的局域形式，因此会移动有效质量参数。}
\end{classroomqa}

\section{每个参数都变成有效参数}

质量图的无限族并不会仅仅因为\(\lambda\)很小就消失。微扰论按耦合的幂次对它们排序，因此相继的项可能越来越小，但仍然有无穷多项。同样的短距离消除也会重整化相互作用本身。

\begin{classroomqa}
\audiencequestion{即使对于质量重整化，不也有无穷多张图吗？\lecturetimestamp{01:00:35}}
\lecturerresponse{确实有。如果\(\lambda\)足够小，微扰级数可能有用，因为更高幂次被相继压低。然而，重整化参数是由整个级数定义的，而不是仅由第一张图定义的。}
\end{classroomqa}

在操作上，一个局域的\(\lambda\phi^4\)相互作用就是一个两粒子进、两粒子出的振幅。如果一个更复杂的四点过程完全发生在一个未解析区域内，那么在低分辨率下它与原来的局域顶点无法区分。因此，这类图被归并到一个有效\(\lambda\)中。它们的主导标度依赖性在四维\(\phi^4\)理论中通常是对数形式的。本讲不计算那个贝塔函数；它用这个例子来说明一般性陈述：
\begin{equation}
\begin{gathered}
\text{parameters measured at low energy}\\
\ne\\
\text{parameters entered at the microscopic cutoff}
\end{gathered}
\end{equation}

四点例子之所以重要，是因为它消除了一种可能的误解。重整化并不是只针对质量发散而施加的特殊修补。拉格朗日量中允许的每一个算符，都可以由具有相同外线和对称性的未解析图来再现。两点过程重整化一个二次系数，四点过程重整化一个四次系数，而真空到真空的过程将重整化常数项。低能拉格朗日量就是一本紧凑的账本，所有这些短距离过程都记录在其中。

\begin{classroomqa}
\audiencequestion{既然\(\Delta\)具有长度的量纲，那么它的对数在量纲上不奇怪吗？还有，黑板上是不是把大写和小写的delta弄反了？\lecturetimestamp{01:01:34}}
\lecturerresponse{对数实际上意味着与另一个标度之比，因此是无量纲的。是的，记号曾被短暂弄反：\(\Delta^\mu\)是积分变量，而\(\delta\)是留在答案中的下距离截断。}
\end{classroomqa}

这是分子的场论对应物。我们解出了看不见的快电子，并把它们的结果放入一个势中。这里我们求和看不见的短距离过程，并把它们的结果放入质量、耦合以及每一个其他允许的局域系数中。一切都发生了移动。

\begin{classroomqa}
\audiencequestion{你能再重复一下，被测量的量是否包含实验看不到的更小标度吗？\lecturetimestamp{01:05:13}}
\lecturerresponse{包含。它们各自的涨落没有被解析，但它们求和后的效应是存在的。所有涉及更短标度的图都被收集到单个有效系数中，低能实验把它称为例如物理质量平方。}
\end{classroomqa}

\begin{classroomqa}
\audiencequestion{那么输入项\(m_0^2\)必须抵消其他贡献吗？\lecturetimestamp{01:06:00}}
\lecturerresponse{如果测得的标量质量远小于各个紫外贡献，那么是的。所需的抵消就是我们接下来要讨论的精细调节问题。}
\end{classroomqa}

\section{希格斯等级问题}

设想截断是普朗克长度，因此有效场论被要求总结从实验室标度一直到量子引力必须介入的标度之间的涨落。
在粒子物理学单位中，
\begin{equation}
\delta\sim M_{\mathrm{Pl}}^{-1}\sim10^{-19}\,\mathrm{GeV}^{-1},
\qquad
\frac{1}{\delta^2}\sim10^{38}\,\mathrm{GeV}^2.
\label{eq:planck-cutoff}
\end{equation}
使用约化还是非约化普朗克标度，对于这个巨大的比较来说无关紧要。

距离等级同样引人注目。普朗克长度比质子半径小约十九个数量级，并且比正在讨论的加速器直接探测的距离小约十六或十七个数量级。正是对这个间隔取平方，才把长度的巨大等级变成了标量质量平方参数中约三十四位的抵消。

对于希格斯场这一标量场，所需的弱电标度是几百GeV量级，而不是\(10^{19}\) GeV。然而，式~\eqref{eq:scalar-mass-series}中的各项却可以自然地具有普朗克尺度的大小。它们的符号尚未被追踪：有些可能为正，有些可能为负，而裸质量平方也可以取任意一种符号。但要留下一个电弱尺度的剩余量，那些巨大的部分必须相消到大约三十四位小数。Susskind 起初说是十七位数字，随后又纠正了自己，因为比较的是质量平方。正是平方之比衡量了这种版本的微调。

\begin{classroomqa}
\audiencequestion{\(\lambda\)是小于还是大于1？\lecturetimestamp{01:09:50}}
\lecturerresponse{假定它略小于1。如果逐次幂增长，微扰论将变得难以处理，而且纯标量理论也无法在任意大的四次耦合下被一致地外推。但\(\lambda\)远没有\(10^{-38}\)那么小；讲座取一个百分数尺度的数值作为参照。}
\end{classroomqa}

\begin{classroomqa}
\audiencequestion{即使\(\lambda<1\)，难道不是仍有许多项因为各自带有大的截断因子而保持重要吗？\lecturetimestamp{01:10:43}}
\lecturerresponse{是的。更高的幂最终会压低级数，但在那发生之前，大量项仍可能作出贡献。它们和裸项都必须共同 conspire，才能留下极其微小的物理答案。}
\end{classroomqa}

回顾一下为什么希格斯质量参数已知处于弱电尺度。按照讲座的约定，将对称破缺势写为
\begin{equation}
V(\phi)=-\frac{\mu^2}{2}\phi^2+\frac{\lambda}{4}\phi^4.
\label{eq:higgs-potential}
\end{equation}
负的二次项系数使原点不稳定，而正的四次项使势稳定。求导，
\begin{equation}
\frac{dV}{d\phi}=-\mu^2\phi+\lambda\phi^3
=\phi(-\mu^2+\lambda\phi^2),
\end{equation}
给出非零极小值
\begin{equation}
\phi_{\min}^2=f^2=\frac{\mu^2}{\lambda}.
\label{eq:higgs-minimum}
\end{equation}
黑板上把四次项示意性地写为\(\lambda\phi^4\)，但引用了式~\eqref{eq:higgs-minimum}；插入约定因子\(1/4\)使这两行一致，而不改变论证。

\(-\mu^2\)的符号很重要，但并不能挽救自然性。正是原点处的负曲率驱动了自发对称破缺。重整化会以大的正、负部分移动该曲率。观测到的破缺真空要求它们的和不仅很小，而且必须具有特定的微小负值，才能把极小值置于弱电尺度。

\(W\)玻色子和\(Z\)玻色子的质量在讲座所用的粗略归一化下确定\(f\)为\(200\,\mathrm{GeV}\)量级。如果\(\lambda\)既不是极大也不是极其微小，那么\(\mu\)必须处于同一大致量级。因此，观测到的弱电尺度系数是那些各自可能大几十个数量级的贡献的剩余。

这被称为希格斯微调问题、等级问题或规范等级问题。等级是指普通弱电尺度与可能的统一尺度或普朗克尺度之间的巨大比值。一个小数本身未必是悖论性的。这里令人不安的是，一个小数由大的、看似无关的部分组成，而这些部分必须以极高精度相消。该问题在1970年代末和1980年代初被认识到，而对这种相消背后机制的探索推动了超对称性以及若干竞争性思想。

无限微扰级数使这个谜题更加尖锐，而不是消除它。百分数尺度的\(\lambda\)使足够高的幂变小，但在这种压低获胜之前，仍可能出现许多项。单圈偶然抵消裸质量是不够的。在纳入其余短距离过程之后，有效弱电尺度必须保持稳定。各项之间的结构性关系将是一种解释；无关的逐位相消只是一种微调。

\section{为什么费米子质量表现更好}

其他标准模型粒子不需要同类的独立微调。对于狄拉克费米子，质量项为
\begin{equation}
\mathcal L_m=-m\bar\psi\psi
=-m\left(\bar\psi_L\psi_R+\bar\psi_R\psi_L\right).
\label{eq:fermion-mass}
\end{equation}
它包含质量，而不是质量的平方，并且耦合左右手征性。因此，局域质量插入把\(L\)变为\(R\)，或把\(R\)变为\(L\)。

\begin{figure}[htbp]
\centering
\includegraphics[width=0.76\textwidth]{lecture_01_figure_04.png}
\caption{费米子质量项以及黑板上改变手征性的线条示意图。由于部分粉笔工作被遮挡，这些图在下面以代数方式重建。}
\lectureframe{01:23:36}
\end{figure}

口头讲座常说是手性或螺旋度。精确的保护最干净地以手征性来表述。对于超相对论性粒子，螺旋度图像提供了预期的直觉，但手征性才是式~\eqref{eq:fermion-mass}中左手场和右手场所携带的量。

规范相互作用保持手征性：
\begin{equation}
L\xrightarrow{\ \gamma\ }L,
\qquad
R\xrightarrow{\ \gamma\ }R.
\label{eq:gauge-preserves-chirality}
\end{equation}
与标量场的汤川相互作用连接相反的手征性：
\begin{equation}
L\xleftrightarrow{\ \phi\ }R.
\label{eq:scalar-flips-chirality}
\end{equation}
首先假设电子以\(m_0=0\)开始。发射并重新吸收任意数量的光子都不能产生质量，因为每个规范顶点都保持手征性。从左开始的序列仍然保持左。因此，纯量子电动力学一致地允许无质量电子：其电磁自能无法从零产生一个可加的质量。

在线图图像中，从左手征电子开始。在第一个光子顶点处它仍然是左手征的；在光子传播并被重新吸收之后，它仍然是左手征的。重复这一构造只会插入更多保持手征性的顶点。外部的两点函数可以被重整化，但它无法获得定义质量所需的左右结构。

标量发射能做到这一点吗？一个标量顶点将\(L\)变为\(R\)，但发射出的标量必须被重新吸收。因此，一个两点自能图包含偶数个标量顶点。第二次翻转逆转了第一次，所以完整的图再次保持外部手征性。无论内部图变得多么复杂，发射和重新吸收将顶点配对，并阻止在起始质量为零时出现可加的\(L\)到\(R\)项。

这是关于翻转次数的宇称论证。一个标量顶点翻转手征性一次；闭合每一条内部标量线都需要第二个端点，因而需要第二次翻转。四个标量顶点翻转四次，依此类推。偶数次使费米子回到其原始手征性。这一结论并不局限于黑板上画出的最简单圈图。

只有当图中已经包含一个质量插入时，非零质量修正才成为可能。例如，光子的发射和重新吸收可以围绕一个\(m_0\bar\psi\psi\)插入。规范顶点贡献两个电荷幂次，但唯一的手征性翻转贡献原始质量：
\begin{equation}
\begin{aligned}
m_e
&=m_0+c_1e^2m_0+c_2e^4m_0+\cdots\\
&=m_0\left(1+c_1e^2+c_2e^4+\cdots\right).
\end{aligned}
\label{eq:multiplicative-fermion-mass}
\end{equation}
如果\(m_0=0\)，结果仍然为零。如果\(m_0\)很小，每个修正都正比于那个小数。我们仍然可以问为什么电子相对于某个基本尺度如此之轻，但我们不再为了得到它而减去大量无关的数。在这个意义上，小的费米子质量在技术上是自然的。

质量插入使允许的序列变得明确。从\(L\)开始。光子发射保持\(L\)；插入\(m_0\)将\(L\)变为\(R\)；光子重新吸收保持\(R\)。该图现在具有质量项的外部手征性，但其值必然包含因子\(m_0\)。这是与式~\eqref{eq:one-loop-scalar-mass}的关键对比，后者的标量修正即使在起始标量质量设为零时仍然非零。

规范玻色子具有类似的保护：规范对称性禁止独立的质量项，因此一个起始无质量的规范玻色子在所有微扰阶都保持无质量。在标准模型中，费米子和弱规范玻色子的质量响应希格斯期望值而产生。控制希格斯尺度，其余轻质量就不需要分别复制同样的灾难性调节。

这确定了超对称性旨在解决的问题：一个理论能否阻止希格斯标量获得巨大的可加质量修正？在适当的超对称设定中，答案是肯定的，因为玻色子和费米子的贡献是相关的，而不是独立的。详细机制从下一讲开始。

\section{另一种精细调节：真空能}

还有一个更为严重的重整化问题，但只有当包含引力时它才变得可观测。在非引力物理中，给每个能量加上一个常数不会改变任何实验。例如，电子质量不是绝对能量；它是向真空添加一个电子所需的额外能量：
\begin{equation}
m_e=E(\text{vacuum}+\text{one electron})-E(\text{vacuum})
\end{equation}
在电子静止系中。如果真空和每个激发态一起移动，它们的差值保持不变。

然而，真空到真空的图确实会移动真空能。在\(\lambda\phi^4\)理论中，将一个四次顶点的四条腿连接成两个短圈。每个重合传播子贡献\(1/\delta^2\)，所以
\begin{equation}
\delta\rho_{\mathrm{vac}}
\sim
c_{\mathrm{vac}}\frac{\lambda}{\delta^4}.
\label{eq:vacuum-energy-divergence}
\end{equation}
更复杂的真空图带有不同的耦合因子和系数，但量纲分析通常为能量密度留下相同的四次截断敏感性。

没有外线正是这一修正的特征。两点图将单粒子态与真空相比较，并改变质量；四点图改变散射；没有入射或出射粒子的图则改变真空到真空的振幅本身。一个四次顶点给出\(\lambda\)，而它的两个短圈各给出\(\delta^{-2}\)。量纲分析随后要求\(\delta^{-4}\)的能量密度标度。

没有引力时，人们可以选择能量的可加零点，并问谁在乎。广义相对论消除了这种自由：应力-能量使时空弯曲，因此均匀的真空能量使空的空间产生引力。它表现为宇宙学常数，并在观测上与暗能量相联系。于是，一个微小的实测真空能量必须经受住形如\(\delta^{-4}\)的贡献。

这就是为什么当引力被开启时问题会发生变化。在普通量子力学中，将每个能级都移动同一个常数不会影响任何跃迁频率。在引力中，这个常数并非惰性的：真空应力张量正比于度规，并使时空弯曲。曾经是任意能量零点的东西，在爱因斯坦方程中变成了一个源。

因此至少存在两个极端的自然性难题。讲座对真空能调节引用的大致是123个十进制位，对希格斯质量调节引用的大致是30--35。确切的数目取决于紫外标度和约定，但这种对比是毋庸置疑的。超对称是对这里讨论的希格斯问题的一种可能回应；一旦超对称破缺，它本身并不能解释观测到的宇宙学常数。\editorialnote{这也是为什么精确超对称的玻色子--费米子相消并不是对观测到的真空能的现象学解决方案：在我们的世界中超对称必须破缺。}

\section{结束课堂提问}

\begin{classroomqa}
\audiencequestion{希格斯修正依赖于截断\(\delta\)，但截断似乎是任意的。它会是普朗克长度、普朗克长度的十倍，还是别的什么？这是否使精细调节变得任意？\lecturetimestamp{01:33:50}}
\lecturerresponse{数值贡献取决于有效理论被信任的最短标度。这种依赖是物理上的记账，而不是逃避。低至标准模型仍然有意义的任何标度的模式，都是其效应必须被包括进来的模式。选择\(\delta\)在新物理被预期出现的距离处---也许是普朗克标度，也许是接近\(10^{16}\,\mathrm{GeV}\)的大统一标度。累积到该标度的贡献已经是该问题的下界。将同一理论延伸到更短的距离只会使二次敏感性更大。我们有义务解释是什么抵消了从实验室到那个紫外边界之间每一个标度的贡献。}
\end{classroomqa}

\begin{classroomqa}
\audiencequestion{什么定义了标量粒子，例子是什么？\lecturetimestamp{01:37:26}}
\lecturerresponse{希格斯是预期的例子。标量在空间旋转下不会变换成不同的分量；等价地说，它的粒子量子具有自旋为零。光子由矢量场描述，而引力场由张量描述。在这次2010讲座时，尚未探测到任何基本自旋为零的粒子。\editorialnote{希格斯玻色子于2012年在CERN被发现，其测量性质与自旋为零一致。}}
\end{classroomqa}

\begin{classroomqa}
\audiencequestion{\(Z\)粒子也是规范玻色子吗？\lecturetimestamp{01:38:31}}
\lecturerresponse{是的。\(Z\)、\(W^+\)和\(W^-\)是自旋为一的规范玻色子，光子也是如此。它们的场方程和相互作用是麦克斯韦式规范描述的非阿贝尔或电弱推广。}
\end{classroomqa}

\begin{classroomqa}
\audiencequestion{为什么玻色子基态贡献在真空能中不抵消费米子贡献？\lecturetimestamp{01:39:10}}
\lecturerresponse{费米子圈和相应的玻色子圈确实可以以相反的符号进入。然而，抵消并非自动的：粒子质量、多重数和耦合必须匹配。即使很小的不匹配也会留下对截断敏感的大残差，而且无限族的图必须逐阶匹配。精确超对称恰好强制了这样的玻色子--费米子关系。一般的非超对称粒子内容则不会，而破缺的超对称会留下残余的真空能问题。}
\end{classroomqa}

\section{展望}

本讲所走的路线是经过深思熟虑的。先通过一个分子计算表明，消去快变量既不是忽略，也不是魔法：快变量的基态能量成为慢变量的势。量子场论在不同波长上重复这一操作。随后量纲分析揭示出哪些有效参数对截断敏感。标量质量会获得正比于\(\delta^{-2}\)的附加项；因此，小的希格斯标度显得不自然地脆弱。费米子质量是乘性重整化的，因为手征性保护了零质量极限，而规范对称性同样保护无质量规范玻色子。真空能量还要更敏感，其标度行为为\(\delta^{-4}\)，而引力使其绝对值成为可观测的。

超对称正是在这一点上进入本课程，它不是作为装饰性的新数学，而是作为一种尝试，使标量部门的行为更像那些受到保护的部门。接下来的任务是理解什么样的对称性能够把玻色子和费米子联系起来，其联系强到足以组织这些抵消。